\documentclass[a4paper]{article}

% Basic packages
\usepackage[fontset=mac]{ctex}
\usepackage{amsmath, amssymb}
\usepackage{graphicx}
\usepackage[margin=2.5cm]{geometry}
\usepackage[most]{tcolorbox}
\usepackage{etoolbox}
\usepackage{listings}
\usepackage{hyperref}
\usepackage{booktabs}
\usepackage{subcaption}
\usepackage{float}
\usepackage{tikz}
\usetikzlibrary{positioning}

% Highlight boxes
\newtcolorbox{knowledgebox}[1]{
    enhanced, colback=blue!5!white, colframe=blue!75!black, colbacktitle=blue!75!black,
    coltitle=white, fonttitle=\bfseries, title=#1, attach boxed title to top left={yshift=-2mm, xshift=2mm},
    boxrule=1pt, sharp corners
}

\newtcolorbox{importantbox}[1]{
    enhanced, colback=yellow!10!white, colframe=yellow!80!black, colbacktitle=yellow!80!black,
    coltitle=black, fonttitle=\bfseries, title=#1, sharp corners
}

\newtcolorbox{warningbox}[1]{
    enhanced, colback=red!5!white, colframe=red!75!black, colbacktitle=red!75!black,
    coltitle=white, fonttitle=\bfseries, title=#1, sharp corners
}

% Listings
\lstset{
    language=Python,
    basicstyle=\ttfamily\small,
    keywordstyle=\color{blue},
    stringstyle=\color{red!60!black},
    commentstyle=\color{green!60!black},
    breaklines=true,
    frame=single,
    numbers=left,
    numberstyle=\tiny\color{gray},
    captionpos=b,
    extendedchars=true,
    inputencoding=utf8
}

% Front-page metadata
\newcommand{\notetitle}{陈天奇：机器学习系统，长期主义，初心，XGBoost，MXNet，TVM，MLC LLM，OctoML｜WhynotTV Podcast \#3}
\newcommand{\noteauthors}{Codex \& 五道口纳什}
\newcommand{\notedate}{\today}
\newcommand{\videochannel}{WhynotTV Podcast}
\newcommand{\videopublishdate}{2025}
\newcommand{\videoduration}{2 小时 40 分钟}
\newcommand{\videourl}{https://www.bilibili.com/video/BV1s6pgzLE3y/}
\newcommand{\videocoverpath}{cover.jpg}

\begin{document}

\begin{titlepage}
\centering
{\Large 播客笔记\par}
\vspace{1.2cm}
{\huge\bfseries \notetitle\par}
\vspace{0.8cm}
{\large \noteauthors\par}
\vspace{0.3cm}
{\large \notedate\par}
\vspace{1.2cm}

\ifx\videocoverpath\empty
{\small 请在生成笔记时填入视频封面图的本地路径。\par}
\else
\includegraphics[width=0.82\textwidth,height=0.45\textheight,keepaspectratio]{\videocoverpath}\par
\fi

\vfill
\begin{tcolorbox}[width=0.9\textwidth, colback=black!2!white, colframe=black!60, sharp corners]
\textbf{视频作者/频道}：\videochannel\par
\textbf{发布日期}：\videopublishdate\par
\textbf{视频时长}：\videoduration\par
\textbf{视频链接}：
\ifx\videourl\empty
未填写
\else
\href{\videourl}{\nolinkurl{\videourl}}
\fi
\end{tcolorbox}
\end{titlepage}

\tableofcontents
\newpage

%% --- 正文内容开始 --- %%

\section{早年经历与成长}

\subsection{童年与计算机的初遇}

陈天奇的计算机启蒙始于少年时期。他在高二时自学编写编译器——这个经历对他后来的技术路线产生了深远影响。编译器是连接高级语言与底层硬件的桥梁，这种"将一种形式翻译为另一种形式"的思维方式，后来贯穿了他从 XGBoost 到 TVM 的整个技术生涯。

\begin{knowledgebox}{编译器的思维方式}
编译器将一个领域特定语言（DSL）的表示翻译为另一个更底层、更可执行的表示。这种"翻译层"的思想，后来在 TVM 中被系统化为"计算与调度分离"（decouple compute from schedule）的核心设计原则——将算法描述（compute）编译为特定硬件的执行方案（schedule）。
\end{knowledgebox}

他在上海交通大学 ACM 班打下系统功底。ACM 班的训练强调算法竞赛与系统编程的结合，这为他日后做高性能机器学习系统提供了稀缺的能力组合：既懂算法，又懂底层系统。

\subsection{ACM 班与系统思维}

ACM 班的经历塑造了陈天奇的系统思维。他提到 ACM 班的导师余荣老师强调"低调做人，高调做事"——这一价值观在他后续的科研生涯中反复体现。与很多追逐热点的研究者不同，陈天奇始终选择"重要但冷门"的方向深耕。

\begin{importantbox}{系统思维与长期主义}
陈天奇反复强调一个核心观点：\textbf{做平庸的问题和做重要的问题花费的时间几乎一样}。因此，选择的问题本身就决定了你能达到的上限。这种"问题选择"的意识，是他从大学阶段就开始养成的思维习惯。
\end{importantbox}

\subsection{本章小结}

陈天奇的早年经历揭示了两个关键线索：编译器思维（从高二自学编译器开始）和系统思维（从 ACM 班训练开始）。这两条线索在他后续的所有项目中都有体现——无论是 XGBoost 对树模型计算图的高效实现，还是 TVM 对深度学习算子的编译优化。

\section{XGBoost：树模型的极致}

\subsection{诞生的背景}

XGBoost（eXtreme Gradient Boosting）诞生于陈天奇在华盛顿大学（UW）攻读博士期间。当时 Kaggle 竞赛正兴，梯度提升树（Gradient Boosting Tree）在结构化数据上表现优异，但现有的实现（如 R 语言的 gbm 包、scikit-learn 的 GBDT）在性能和可扩展性上存在明显短板。

\begin{figure}[H]
\centering
\begin{tikzpicture}[node distance=2.5cm, auto]
  \tikzstyle{block}=[rectangle, draw, minimum width=3.2cm, minimum height=1cm,
                     align=center, font=\small, rounded corners=2pt]
  \tikzstyle{arrow}=[thick,->,>=stealth]

  \node [block, fill=blue!10] (data)    {训练数据};
  \node [block, fill=green!10, right of=data] (boost)   {梯度提升迭代};
  \node [block, fill=yellow!10, right of=boost] (tree)   {CART 决策树};
  \node [block, fill=red!10, right of=tree] (pred)   {集成预测};

  \draw [arrow] (data) -- (boost);
  \draw [arrow] (boost) -- (tree);
  \draw [arrow] (tree) -- (pred);
  \draw [arrow] (pred.south) to [out=-30, in=-150] (boost.south);
\end{tikzpicture}
\caption{XGBoost 的核心训练流程：每轮迭代拟合前一轮的残差，所有树的预测加权求和}
\end{figure}

\subsection{核心设计哲学}

陈天奇对 XGBoost 的设计哲学可以概括为三个维度：

\begin{enumerate}
\item \textbf{模型的容量不应受限}：树模型的一个关键特性是"可任意增长"（arbitrary growth）——随着数据量增加，树的深度和数量都可以相应扩展。这与固定参数数量的线性模型形成鲜明对比。

\item \textbf{系统级优化}：XGBoost 不只是算法实现，更是系统工程。它引入了列块（column block）数据结构来加速分裂点搜索、利用缓存感知（cache-aware）的内存访问模式、支持分布式训练和核外计算（out-of-core）。

\item \textbf{可扩展的正则化}：在传统 GBDT 的目标函数中显式加入正则化项（树的叶子节点权重和叶子数量），使其比传统方法更不容易过拟合。
\end{enumerate}

\begin{importantbox}{XGBoost 目标函数}
XGBoost 在第 $t$ 轮迭代时的目标函数为：
$$
\mathcal{L}^{(t)} = \sum_{i=1}^{n} l(y_i, \hat{y}_i^{(t-1)} + f_t(x_i)) + \Omega(f_t)
$$
其中 $\Omega(f_t) = \gamma T + \frac{1}{2}\lambda \sum_{j=1}^{T} w_j^2$ 是正则化项。
\begin{itemize}
\item $l$：损失函数（如平方误差、对数损失）
\item $\hat{y}_i^{(t-1)}$：前 $t-1$ 轮的累积预测
\item $f_t$：第 $t$ 棵树的输出
\item $T$：叶子节点数，$w_j$：第 $j$ 个叶子的权重
\item $\gamma$、$\lambda$：正则化超参数
\end{itemize}
\end{importantbox}

\subsection{Kaggle 时代的统治力}

陈天奇回忆，XGBoost 在 2015 年前后的 Kaggle 竞赛中几乎成为"标配"工具。29 个 Kaggle 冠军解决方案中，有 17 个使用了 XGBoost。但陈天奇对此保持清醒——他提到，"如果我们继续去优化 ResNet，最后会输得很惨"。这一判断直接推动了 MXNet 的诞生。

\subsection{本章小结}

XGBoost 的成功证明了陈天奇的一个核心信念：\textbf{系统工程的深度决定了算法能走多远}。XGBoost 不仅是算法的胜利，更是系统工程的胜利——高效的数据结构、精细的内存管理、可扩展的分布式设计，这些才是它在 Kaggle 中碾压对手的真正原因。

\section{MXNet：深度学习框架之争}

\subsection{从 XGBoost 到 MXNet}

陈天奇回忆，MXNet 的起点是对"下一步该做什么"的思考。在 XGBoost 获得成功后，他意识到树模型在非结构化数据（图像、语音、文本）上的能力有限，深度学习才是未来的大方向。

MXNet 的早期团队成员包括李沐（Mu Li）、Alex Smola 等，他们在 UW 和 CMU 联合推进。陈天奇描述当时的氛围："我们刚开始做 MXNet 的时候，完全不会担心 TensorFlow 能对我们产生任何威胁。"

\begin{figure}[H]
\centering
\begin{tikzpicture}[node distance=0.5cm, auto]
  \tikzstyle{layer}=[rectangle, draw, minimum width=10cm, minimum height=0.7cm,
                     align=center, font=\small, rounded corners=1pt]
  \node [layer, fill=blue!10] (api)    {Python / Scala / R / Julia / C++ 多语言前端};
  \node [layer, fill=green!10, below of=api] (graph)  {符号式（Symbolic）+ 命令式（Imperative）混合编程};
  \node [layer, fill=yellow!10, below of=graph] (engine) {依赖引擎（Dependency Engine）— 自动并行与调度};
  \node [layer, fill=red!10, below of=engine] (backend) {后端适配层：CPU / GPU / 分布式};
\end{tikzpicture}
\caption{MXNet 架构概览：混合式编程模型 + 依赖引擎 + 多后端适配}
\end{figure}

\subsection{Python-First 的先见之明}

在 MXNet 诞生的年代（2014--2015 年），主流的深度学习框架生态是：

\begin{itemize}
\item Caffe：C++ 核心，Python 绑定，但以配置文件（prototxt）驱动为主
\item Torch7：Lua 语言，生态较小
\item Theano：Python，但编译慢，调试困难
\item TensorFlow：刚发布，静态图为主
\end{itemize}

陈天奇团队做了一个在当时看来不显然的决定：\textbf{Python First}。他们同时支持符号式编程和命令式编程（类似今天的 PyTorch eager mode），这一设计后来被证明是正确的方向。

\begin{knowledgebox}{符号式 vs 命令式编程}
\begin{itemize}
\item \textbf{符号式（Symbolic）}：先定义计算图，后执行。优势是图优化、内存复用、跨设备部署。典型代表：TensorFlow 1.x、Theano。
\item \textbf{命令式（Imperative）}：逐行执行，即写即得。优势是调试方便、控制流自然。典型代表：PyTorch、NumPy。
\item MXNet 的创新：\textbf{同时提供两种接口}，用户可以在同一个模型中混用——这是后来 JAX 等框架也在探索的方向。
\end{itemize}
\end{knowledgebox}

\subsection{框架战争的教训}

当被问到对框架战争的看法时，陈天奇的反思格外坦诚。他认为 MXNet 输给 PyTorch/TensorFlow 的核心原因是：

\begin{enumerate}
\item \textbf{生态建设不足}：模型 zoo、预训练权重、教程社区的建设比对手晚了太多
\item \textbf{团队资源差距}：Google（TensorFlow）和 Facebook（PyTorch）可以投入数倍甚至数十倍的工程资源
\item \textbf{时机与路径依赖}：当学术界/工业界已经开始使用某个框架，切换成本极高
\end{enumerate}

\begin{warningbox}{框架战争的深层教训}
陈天奇指出一个容易忽略的事实：框架的胜负不取决于技术优劣。MXNet 在技术架构上有很多先进设计（混合编程、依赖引擎、轻量级），但"技术先进"不是框架取胜的充分条件。生态、社区、资源投入、时机，这些因素在系统工程中往往比技术本身更重要。
\end{warningbox}

\subsection{本章小结}

MXNet 的经历让陈天奇更深刻地理解了"现实"——如何通过现实的方式去实现理想。他坦言这段经历让他的做事方式变得更加务实。从 MXNet 的教训中，他提炼出了 TVM 的方向：不再做框架，而是做框架之下的编译基础设施。

\section{TVM：深度学习编译的开创}

\subsection{问题起源}

MXNet 之后，陈天奇开始思考一个更深层的问题：\textbf{如何让深度学习模型高效运行在各种不同的硬件上？} 这个问题在当时几乎无人问津——因为当时的共识是"找几家大厂的工程师手工为每个硬件写 kernel 就好了"。

但陈天奇从编译器的视角看到了一种不同的可能性：深度学习算子本质上是一种计算模式（compute pattern），如果能将计算描述与硬件调度分离，就可以自动生成针对不同硬件的优化代码。

\begin{figure}[H]
\centering
\begin{tikzpicture}[node distance=0.6cm, auto]
  \tikzstyle{layer}=[rectangle, draw, minimum width=10cm, minimum height=0.8cm,
                     align=center, font=\small, rounded corners=1pt]
  \node [layer, fill=blue!10] (model)    {深度学习模型（Relay IR）};
  \node [layer, fill=green!10, below of=model] (graph)  {图级优化：算子融合、常量折叠、内存规划};
  \node [layer, fill=yellow!10, below of=graph] (tir)   {张量级 IR（TIR）：计算与调度分离};
  \node [layer, fill=orange!10, below of=tir] (schedule) {自动调度 / AutoTVM / AutoScheduler};
  \node [layer, fill=red!10, below of=schedule] (codegen) {代码生成：LLVM / CUDA / OpenCL / Metal / Vulkan};
\end{tikzpicture}
\caption{TVM 编译栈：从模型到硬件指令的多层编译流水线}
\end{figure}

\subsection{计算与调度分离}

TVM 最核心的设计理念是 \textbf{Decouple Compute from Schedule}：

\begin{itemize}
\item \textbf{Compute}：描述"做什么"——数学运算本身，与硬件无关，如矩阵乘法的数学定义
\item \textbf{Schedule}：描述"怎么做"——在特定硬件上的优化策略，如 tiling、vectorization、parallelization、cache hierarchy
\end{itemize}

\begin{importantbox}{TVM 核心抽象}
一个 TVM 算子由两部分组成：
$$
\text{算子} = \text{Compute（计算描述）} + \text{Schedule（调度策略）}
$$
\begin{itemize}
\item Compute 是平台无关的数学定义，只需写一次
\item Schedule 是平台相关的优化方案，通过自动搜索（AutoTVM）或手写模板生成
\item 同一套 Compute 可以通过不同的 Schedule 生成针对 CPU、GPU、NPU 等不同后端的代码
\end{itemize}
\end{importantbox}

\subsection{VTA：可编程 NPU 的前瞻设计}

TVM 的一个子项目 VTA（Versatile Tensor Accelerator）展示了对 NPU 架构的前瞻思考。陈天奇指出，当时大多数 NPU 是"完全硬连线的"（fully hardwired）——不可编程，只能做固定功能。VTA 从编译器的角度出发，设计了一套\textbf{可编程的微指令集}，使得同一个加速器可以执行多种深度学习算子。

\begin{knowledgebox}{VTA 的设计洞察}
VTA 的关键洞察：将矩阵运算分解为更小的步骤（小矩阵乘法和内存读写），用一套统一的指令集编程，通过指令重叠（overlap）实现高效执行。这种设计使得 NPU 不再是固定的硬件功能块，而是一个可编程的张量处理器——在 ASIC 的效率和通用处理器的灵活性之间找到了平衡点。
\end{knowledgebox}

\subsection{本章小结}

TVM 代表了陈天奇从"做框架"到"做编译基础设施"的战略转型。与其在框架层面与 Google/Facebook 竞争，不如在更底层的编译优化层面建立自己的壁垒。这种"向下沉"的策略后来被证明非常成功——TVM 成为 Apache 顶级项目，被亚马逊、华为、高通等公司广泛采用。

\section{MLC LLM：让大模型无处不在}

\subsection{从 TVM 到 MLC}

MLC（Machine Learning Compilation）LLM 是 TVM 思路在 LLM 时代的延续和升级。陈天奇回忆，在 2022 年底 ChatGPT 发布后，他很快意识到一个新的问题：\textbf{如何让大语言模型运行在每个人的设备上，而不只依赖云端 API？}

\begin{figure}[H]
\centering
\begin{tikzpicture}[node distance=0.6cm, auto]
  \tikzstyle{layer}=[rectangle, draw, minimum width=10cm, minimum height=0.8cm,
                     align=center, font=\small, rounded corners=1pt]
  \node [layer, fill=blue!10] (model)    {模型格式：HuggingFace / PyTorch / GGUF / …};
  \node [layer, fill=green!10, below of=model] (quantize)  {量化压缩：int4 / int3 / 混合精度};
  \node [layer, fill=yellow!10, below of=quantize] (compile) {MLC 编译：TVM Unity IR → 平台代码};
  \node [layer, fill=orange!10, below of=compile] (runtime) {运行时：CPU / GPU / Metal / Vulkan / WebGPU};
  \node [layer, fill=red!10, below of=runtime] (deploy) {部署目标：iPhone / Android / 笔记本 / 浏览器};
\end{tikzpicture}
\caption{MLC LLM 流水线：从预训练模型到全平台部署}
\end{figure}

\subsection{Personal Computer 的类比}

陈天奇用了一个精妙的类比来解释 MLC LLM 的愿景：

\begin{importantbox}{Mainframe → Personal Computer 的历史类比}
历史上，计算机经历了从大型机（Mainframe）到个人电脑（Personal Computer）的转变。今天的大语言模型正处在一个类似的转折点：
\begin{itemize}
\item \textbf{Mainframe 时代}：所有计算集中在云端，用户通过终端访问。对应今天的 ChatGPT API 模式。
\item \textbf{PC 时代}：每个人有自己的计算设备，本地运行软件。对应未来每个人设备上运行本地 LLM。
\item \textbf{关键前提}：本地 LLM 的能力需要足够强。MLC 的工作就是让这个"足够强"更快到来。
\end{itemize}
\end{importantbox}

\subsection{技术挑战与路线}

MLC LLM 面临的核心技术挑战：

\begin{enumerate}
\item \textbf{量化}：大模型参数量巨大（7B--70B+），需要通过 int4/int3 量化在精度和速度之间找到平衡
\item \textbf{跨平台}：iPhone（Metal）、Android（OpenCL/Vulkan）、笔记本（CUDA/CPU）、浏览器（WebGPU）各平台差异巨大
\item \textbf{内存受限}：移动设备内存通常只有 4--8GB，需要精细的内存规划和 KV-cache 管理
\end{enumerate}

\subsection{本章小结}

MLC LLM 将 TVM 的编译思想从"通用深度学习模型"聚焦到了"大语言模型"，并且在目标平台上从"云端/服务器"扩展到"个人设备/浏览器"。陈天奇将 MLC LLM 定位为"让机器学习系统工程像写神经网络一样简单"——这是他从 XGBoost 时代延续至今的愿景。


\section{机器学习系统：作为一个学科}

\subsection{为什么要定义"机器学习系统"}

在访谈中，陈天奇花了相当长的篇幅阐述他对"机器学习系统"（Machine Learning Systems）作为一个学科的看法。他指出，机器学习的历史通常被讲述为"算法进步史"（AlexNet → ResNet → Transformer → GPT），但\textbf{系统层面的进步同样关键，甚至更关键}——没有高效的训练和推理系统，这些算法的实际价值将大打折扣。

\begin{importantbox}{机器学习系统的定义}
机器学习系统是研究\textbf{如何高效实现、训练、部署机器学习模型}的学科，涵盖：
\begin{itemize}
\item \textbf{训练系统}：分布式训练、混合精度、梯度压缩、流水线并行
\item \textbf{推理系统}：模型量化、算子融合、内存优化、批处理调度
\item \textbf{编译系统}：计算图优化、算子代码生成、硬件后端适配
\item \textbf{部署系统}：模型服务、边缘端推理、跨平台适配
\end{itemize}
\end{importantbox}

\subsection{从"调参"到"系统工程"}

陈天奇观察到，当前 ML 系统工程的一个瓶颈是门槛仍然很高。"你写一个神经网络已经非常简单了，在 Python 里面写几行就可以做到。但做机器学习系统工程开发还远远没有那么简单。"他希望能让 ML 系统工程也变得像写 PyTorch 代码一样简单——这正是 MLC 项目的使命。

\begin{figure}[H]
\centering
\begin{tikzpicture}[node distance=2.0cm, auto]
  \tikzstyle{block}=[rectangle, draw, minimum width=3cm, minimum height=0.9cm,
                     align=center, font=\small, rounded corners=2pt]
  \tikzstyle{arrow}=[thick,->,>=stealth]

  \node [block, fill=blue!10] (algo)    {算法创新\\（论文驱动）};
  \node [block, fill=green!10, right of=algo, xshift=1.5cm] (sys)   {系统创新\\（工程驱动）};
  \node [block, fill=yellow!10, below of=algo, yshift=-0.5cm] (data)  {数据/算力};
  \node [block, fill=red!10, below of=sys, yshift=-0.5cm] (deploy) {部署/落地};

  \draw [arrow] (algo) -- (sys) node[midway, above, font=\small] {提出需求};
  \draw [arrow] (sys) -- (deploy) node[midway, right, font=\small] {使能};
  \draw [arrow] (data) -- (algo);
  \draw [arrow] (data) -- (sys);
  \draw [arrow, dashed] (sys.west) to [out=180, in=270] (algo.south) node[near start, left, font=\small] {反馈约束};
\end{tikzpicture}
\caption{算法与系统的共生关系：系统使能算法，算法提出系统需求，两者在约束中互相推动}
\end{figure}

\subsection{边缘 AI 的想象空间}

陈天奇对边缘 AI（Edge AI）的前景非常乐观。他将当前的 LLM 生态类比为计算机发展史上的 Mainframe 时代——所有智能集中在云端数据中心。但正如 Personal Computer 最终改变了计算的面貌，他相信\textbf{本地 LLM}也会改变 AI 的使用方式。

\begin{knowledgebox}{为什么需要本地 LLM}
\begin{itemize}
\item \textbf{隐私}：敏感数据（个人笔记、医疗记录、金融信息）不应离开用户设备
\item \textbf{延迟}：云端推理的网络延迟（100--500ms）对实时交互（如语音助手）是致命问题
\item \textbf{离线可用}：飞机上、地下室、野外——网络不是随时可用的
\item \textbf{成本}：云端 GPU 推理成本高昂，本地推理边际成本接近零
\item \textbf{个性化}：本地模型可以与用户的个人数据深度绑定，提供定制化体验
\end{itemize}
\end{knowledgebox}

\subsection{本章小结}

"机器学习系统"的学科界定是陈天奇对领域的一个重要贡献。他不仅做了 XGBoost、MXNet、TVM、MLC 这几个具体项目，更通过这些项目逐步定义和充实了"机器学习系统"这个学科的内涵——从"调参工具"到"系统工程"，从"少数人的手艺"到"可复用的基础设施"。

\section{OctoML：创业的试炼}

\subsection{从 Apache TVM 到 OctoML}

OctoML 是陈天奇与 TVM 核心团队在 2019 年创立的公司，目标是将 TVM 技术商业化。陈天奇坦率地分享了这段经历的种种挑战和认知转变。

\begin{importantbox}{OctoML 商业模式的三次迭代}
\begin{enumerate}
\item \textbf{阶段一（2019--2021）}：做 TVM 的企业级工具/平台，目标客户是硬件厂商和 ML 部署团队。\\问题：市场太小，TVM 的开源版本已经足够好。
\item \textbf{阶段二（2022--2023）}：转型做 ML 部署 SaaS 服务——类似 Fireworks、Together AI 的商业模式，帮客户部署和优化模型。\\问题：竞争激烈，大模型 API 服务的差异化不够明显。
\item \textbf{阶段三（2023--至今）}：聚焦特定垂直领域的端到端优化方案。\\方向：让"部署"这件事本身变得更简单、更标准化。
\end{enumerate}
\end{importantbox}

\subsection{创业 vs 科研}

陈天奇分享了创业与科研的深刻比较：

\begin{itemize}
\item \textbf{科研}：追求"从 0 到 1"的突破，失败是可接受的结果，探索本身就有价值
\item \textbf{创业}：追求"Product-Market Fit"，必须找到真正的付费客户，商业逻辑比技术先进性更重要
\item \textbf{互补}：创业让他"更现实"——知道如何通过现实的方式去实现理想
\end{itemize}

\subsection{本章小结}

OctoML 的经历让陈天奇从纯粹的研究者变成了一个更全面的 builder。他形容这段经历"让我更现实，我会知道怎么样能够通过现实的方式去更好地实现理想"。这种"理想主义 + 现实主义"的双重思维，是他与其他纯学术研究者最显著的区别之一。

\section{CMU 教职与指导哲学}

\subsection{为什么选择教职}

陈天奇在博士毕业时就决定走教职路线。他的理由有三个：

\begin{enumerate}
\item \textbf{喜欢科研的不确定性}：做没有人知道答案的问题，这种探索感是不可替代的
\item \textbf{喜欢和同学们一起创造}：他提到"和同学们一起去解决新的东西非常有趣"
\item \textbf{喜欢教学}：把自己的知识和经验传递下去，帮助下一代研究者成长
\end{enumerate}

\subsection{Carlos Guestrin 的影响}

陈天奇多次提到博士导师 Carlos Guestrin 对他的深远影响。Guestrin 不仅是 XGBoost 项目的共同推动者，更在科研方法论上给陈天奇留下了深刻的印记：

\begin{knowledgebox}{Guestrin 的科研方法论}
\begin{itemize}
\item \textbf{做有影响力的问题}：不要满足于增量改进，要寻找"能改变游戏规则"的问题
\item \textbf{论文要 accessible}：好的研究应该让读者容易理解，而不是用复杂公式吓退别人
\item \textbf{从 Research 到 Adoption}：一个好的研究项目应该能够从 paper 阶段推动到真正被人使用的阶段
\end{itemize}
\end{knowledgebox}

\subsection{作为导师的风格}

陈天奇描述自己的指导风格：延续 Guestrin 的很多理念，但同时加入了自己的特色。他强调：

\begin{itemize}
\item \textbf{团队合作}：机器学习系统是一个需要团队协作的方向，单打独斗很难走远
\item \textbf{做生的、做有影响力的工作}：鼓励学生选择"难但重要"的问题，而不是"容易发论文"的问题
\item \textbf{跨领域视野}：希望学生能够同时理解算法、系统、硬件三个层面
\end{itemize}

\subsection{本章小结}

陈天奇的教职身份（CMU Assistant Professor）与他的创业者身份（OctoML co-founder）形成了一种有趣的张力。他自己描述为"80\% 教授，20\% 创业者"——这种双重身份让他能够在学术深度和产业落地之间找到平衡。

\section{长期主义与价值观}

\subsection{被问到"想对十年前的自己说什么"时的沉默}

这期播客最打动人心的片段之一，是主持人对陈天奇的一个提问：\textbf{"你会对十年前、二十年前的陈天奇说些什么？"} 陈天奇沉默了近 20 秒后回答：

\begin{quote}
"可能要反过来——我需要过去的我对现在的我，现在的我对未来的我说：记住自己对自己的承诺，坚持自己的理想，往下走下去。"
\end{quote}

这个回答揭示了陈天奇价值观的核心：\textbf{不是"现在的我"去纠正"过去的我"，而是"过去的初心"来提醒"现在的我"不要迷失}。

\begin{importantbox}{长期主义的实践}
陈天奇对长期主义的理解不是抽象的口号，而是非常具体的操作原则：
\begin{enumerate}
\item \textbf{问题选择}：做重要的问题和平庸的问题花费的时间几乎一样——选择本身就决定了上限
\item \textbf{接受失败}：做难的问题必然意味着不确定性和失败，关键在于"be willing to take the risk"
\item \textbf{Reinvent yourself}：从 XGBoost 到 MXNet 到 TVM 到 MLC，每一次都是一次自我革新
\item \textbf{勇气来自过去的自己}：当你犹豫不决时，回想自己最初为什么出发——那份初心会给你继续的勇气
\end{enumerate}
\end{importantbox}

\subsection{为什么"现在保持初心更难"}

当陈天奇说"现在保持初心更难"时，他提到了一个悖论：

\begin{itemize}
\item \textbf{资源多了}：现在的陈天奇能调动的资源是十年前的"不是一个数量级"
\item \textbf{选择多了}：但选择越多，反而越容易迷失——"你会开始担心"
\item \textbf{对比效应}：当你看到别人走得更快、做得更"热"时，会动摇自己对冷门方向的坚持
\end{itemize}

\begin{warningbox}{选择的悖论}
陈天奇指出一个常见的陷阱：当资源充足时，反而更容易做平庸的决策——因为"安全"的路径看起来更加诱人。他说："我们刚开始做 MXNet 的时候，完全不会担心 TensorFlow 能对我们产生什么威胁。"那种"一无所有所以无所畏惧"的状态，反而是创新的最佳土壤。
\end{warningbox}

\subsection{给年轻研究者的建议}

陈天奇对年轻研究者（特别是 PhD 学生）的建议可以浓缩为三条：

\begin{enumerate}
\item \textbf{选择自己真正有 passion 的问题}：不要跟风，不要为了发论文而选择"安全"的方向
\item \textbf{要有系统级的视野}：既能看到具体的技术问题，也能理解问题在整个系统中的位置
\item \textbf{不要害怕失败}：他引用自己早期的失败经历——手搓 RBM 做 ImageNet classification 失败，ACM 竞赛失败——来说明失败本身就是科研的常态
\end{enumerate}

\subsection{本章小结}

陈天奇的价值观可以概括为：\textbf{长期主义 × 敢失败 × 系统思维}。这是他在 20 年机器学习系统工程师生涯中提炼出的底层哲学，也是他所有技术决策的最终依据。

\section{总结与延伸}

\subsection{技术路线的演进逻辑}

回顾陈天奇 20 年的技术路线，可以看到一条清晰的演进逻辑：

$$
\text{XGBoost（树模型极致）} \to \text{MXNet（框架层）} \to \text{TVM（编译层）} \to \text{MLC LLM（全栈部署）}
$$

这条路线中有几个关键的战略转折：

\begin{enumerate}
\item \textbf{从树模型到深度学习}（2014--2015）：意识到树模型的上限，主动转向深度学习框架
\item \textbf{从框架到编译}（2016--2017）：MXNet 在框架战中落败后，"向下沉"到编译基础设施
\item \textbf{从云端到边缘}（2022--至今）：MLC LLM 将焦点从服务器端推理转向个人设备部署
\end{enumerate}

\begin{figure}[H]
\centering
\begin{tikzpicture}[node distance=1.2cm, auto]
  \tikzstyle{phase}=[rectangle, draw, minimum width=4cm, minimum height=0.9cm,
                     align=center, font=\small, rounded corners=2pt, fill=blue!5]
  \tikzstyle{arrow}=[thick,->,>=stealth]

  \node [phase] (xgb)     {XGBoost\\2014};
  \node [phase, right of=xgb, xshift=2cm] (mxnet)  {MXNet\\2015--2017};
  \node [phase, right of=mxnet, xshift=2cm] (tvm)    {TVM\\2017--2022};
  \node [phase, right of=tvm, xshift=2cm] (mlc)    {MLC LLM\\2022--};

  \node [font=\footnotesize, below of=xgb, yshift=0.2cm] {树模型系统};
  \node [font=\footnotesize, below of=mxnet, yshift=0.2cm] {框架层};
  \node [font=\footnotesize, below of=tvm, yshift=0.2cm] {编译层};
  \node [font=\footnotesize, below of=mlc, yshift=0.2cm] {全栈部署};

  \draw [arrow] (xgb) -- (mxnet) node[midway, above, font=\small] {向上};
  \draw [arrow] (mxnet) -- (tvm) node[midway, above, font=\small] {向下};
  \draw [arrow] (tvm) -- (mlc) node[midway, above, font=\small] {铺开};
\end{tikzpicture}
\caption{陈天奇 20 年技术路线的演进：向上（框架）→ 向下（编译）→ 铺开（全平台）}
\end{figure}

\subsection{核心主张}

从 2 小时 40 分钟的深度访谈中，可以蒸馏出陈天奇的几个核心主张：

\begin{enumerate}
\item \textbf{系统价值被低估}：ML 的进步不只是算法的故事，系统工程的深度同样决定了模型能走多远。XGBoost 证明了这一点，TVM 也证明了这一点。
\item \textbf{选择本身就是最重要的决策}：做平庸的问题和做重要的问题花费的时间几乎一样。问题选择的眼界，来自长期积累的系统视野。
\item \textbf{在不确定性中找到确定性}：科研和创业都是高不确定性活动，"be willing to take the risk"不是勇气问题，而是认知问题——理解失败是过程的一部分。
\item \textbf{不忘初心不是道德要求，是战略要求}：从过去的自己那里汲取勇气，不是因为"坚持"本身是美德，而是因为只有初心能告诉你什么是真正值得做的。
\item \textbf{现实是理想的实现方式}：OctoML 的创业经历教会陈天奇"通过现实的方式去更好地实现理想"——理想主义不排斥现实主义，相反，现实主义是理想主义的执行手段。
\end{enumerate}

\subsection{开放问题}

陈天奇在访谈中留下了几个值得思考的开放问题：

\begin{itemize}
\item \textbf{本地 LLM 的能力何时能达到"足够好"的水平？} 这决定了边缘 AI 的时间表
\item \textbf{机器学习系统能否真正成为一个独立的学科？} 目前它仍介于 CS Systems 和 ML 之间，边界模糊
\item \textbf{当下一个范式转换来临时，什么会成为新的瓶颈？} 陈天奇的回答是"继续去解决让更少的人用更少的 engineering resource 去解决 ML 系统的问题"
\item \textbf{如何在资源充裕时保持"敢失败"的心态？} 这是陈天奇自己也在思考的问题
\end{itemize}

\subsection{拓展阅读}

\begin{itemize}
\item 陈天奇个人主页：\href{https://tqchen.com/}{tqchen.com}
\item 陈天奇知乎文章「机器学习科研的十年」：\href{https://zhuanlan.zhihu.com/p/74249758}{zhuanlan.zhihu.com/p/74249758}
\item WhynotTV 播客主页：\href{https://tairanhe.com/podcast\_cn}{tairanhe.com/podcast\_cn}
\item Apache TVM 项目：\href{https://tvm.apache.org/}{tvm.apache.org}
\item MLC LLM 项目：\href{https://github.com/mlc-ai/mlc-llm}{github.com/mlc-ai/mlc-llm}
\end{itemize}

\end{document}
